\section{Audit Instrument and Conditional Guarantees}
\label{sec:system}
\label{sec:graph_model}
\label{sec:verification_methods}

% Pipeline figure moved to the supplement to meet the page limit.

\paragraph{Unified graph model.}
The analyzer represents a workflow as a directed graph $G=(V,E,v_0,V_f)$ with
entry $v_0$ and exit set $V_f$. Each node carries a kind $\kappa_V(v)$ over
\textsc{entry}, \textsc{exit}, \textsc{tool}, \textsc{llm}, \textsc{router},
\textsc{human}, \textsc{passthrough}, \textsc{subgraph} and, for tool nodes,
bound tool names; each edge is \textsc{direct}, \textsc{conditional},
\textsc{parallel}, or \textsc{loop}. Every node and edge also carries extraction
\emph{provenance}---an origin (runtime object, explicit AST, inferred,
synthesized, unknown), a confidence (exact, may, heuristic), and a source
span---so structure read from a declaration is distinguishable from structure the
extractor guessed, and witness traces inherit the weakest confidence they
traverse. Per-framework extractors bridge four representations into this model
by either of two paths: a runtime extractor reading the compiled graph object,
or a static AST fallback for large-scale mining, where installing every
repository's dependencies is infeasible.

\paragraph{Event attributes.}
\label{par:ontology}
Nodes may additionally carry non-exclusive effects (read, write, execute,
communicate, financial, delete) and capabilities. The mined extractors do not
populate these fields from node bodies; corpus construction assigns them from
declared tool names. This limitation is central to the audit rather than hidden
by the representation.

\paragraph{Structural checks.}
Six checks run over the topology, each in $O(|V|+|E|)$ and each emitting a
\emph{witness trace}~\citep{clarke1999modelchecking} on failure: \emph{exit
reachability} ($V_f\subseteq\mathrm{Reach}(v_0)$), \emph{reverse reachability}
(reachable nodes with no path to an exit), \emph{dead-end detection} (non-exit
nodes with no successor), \emph{router shape} (router outgoing edges must be
conditional), \emph{human presence} (no \textsc{human} node---a deliberately
blunt policy check we evaluate against a risk-aware variant), and \emph{tool
declaration} (tool nodes with an empty tool set). A seventh, path-sensitive
check---\emph{human-gate coverage}---flags a workflow only if a \emph{sensitive}
tool is reachable from $v_0$ without passing a \textsc{human} node
(Section~\ref{sec:realworld}).

\paragraph{Temporal policies with LTL$_f$ semantics.}
It also supports a policy DSL over finite traces
(LTL$_f$~\citep{degiacomo2013ltlf}): an execution is a finite event sequence, and
each of seven DSL forms denotes an LTL$_f$ formula---forbidden ($\mathbf{G}\,\neg
a$), response ($\mathbf{G}(a \rightarrow \mathbf{F}\,b)$), strong until, bounded
response ($\mathbf{G}(a \rightarrow \mathbf{F}_{\leq k}\,b)$), response chain,
conjunction, disjunction (grammar in the supplement). Formulas compile
compositionally to a DFA by formula progression (2--3 states for base patterns,
$k{+}2$ for bounded response), with two disjoint violation mechanisms:
\emph{violation states} are exactly the doomed states, so a bad prefix is
flagged at the earliest event making every extension violating; and a pending
$\mathbf{F}$-obligation violates at termination. The compiler is differentially
tested against an independent reference evaluator, exhaustively to length six
for two-atom formulas. Runtime monitors evaluate events in $O(1)$; obligations
open when a run aborts are reported inconclusive~\citep{bauer2011ltl3}.

\paragraph{Static temporal verification.}
Policies are checked statically via a graph\,$\times$\,DFA product: BFS from
$(v_0,q_0)$ advances the DFA at each node and follows graph edges, modeling
executions as finite \emph{maximal} paths. Because static extraction cannot prove
which of a tool node's declared tools will be invoked, or whether any will be,
the product is \emph{conservative}: at each tool node it closes the DFA state
over every finite invocation sequence drawn from the declared tools, including
the empty one---a false alarm is acceptable, a false proof is not. The checker
returns \emph{may-violate} (some resolution reaches a bad prefix or a pending
obligation at termination), \emph{safe} (none does), or \emph{inconclusive} (a
reachable product cycle postpones an obligation forever, or no termination point
is reachable; an engineering-scope choice, not undecidability). Explored states
are bounded by $|V|\cdot|Q|$.

\paragraph{Soundness requires an event-complete abstraction.}
\label{thm:soundness}
Soundness is a claim about \emph{event} traces, not graph paths. Let
$\mathrm{Traces}_{\mathrm{ev}}(P)$ be the concrete finite terminating event traces
a workflow $P$ can emit, and let $\lambda$ label each event, inducing a graph
event-language $\mathcal{L}(G,\lambda)$. If $\mathrm{Traces}_{\mathrm{ev}}(P)\subseteq
\mathcal{L}(G,\lambda)$ (\emph{event-trace conservatism}) and every trace in
$\mathcal{L}(G,\lambda)$ is accepted by the policy DFA, then every concrete finite
terminating execution of $P$ satisfies the policy---the guarantee we exploit for
monitor pruning (Section~\ref{sec:results}); the proof is in the supplement.
Graph-\emph{path} containment alone does \emph{not} suffice: a node may issue tool
calls the graph never records. Because provenance cannot establish that node
bodies emit no unmodeled events, the checker returns \texttt{safe} only when the
caller explicitly asserts event-trace conservatism for an authored abstraction.
Without that assertion, a would-be safe result is downgraded to
\texttt{inconclusive} with reason \texttt{uncertified\_extraction};
\texttt{may\_violate} findings remain visible.

\paragraph{What the gate does and does not certify.}
Exact provenance describes \emph{how} an element was obtained, not whether the
abstraction covers all runtime events. An edge can be read exactly while a tool
call inside its node body is absent. We therefore use provenance to qualify
witnesses, never to discharge event-trace conservatism. The product is sound
\emph{relative to an authored abstraction} whose event coverage the caller
asserts; on every mined graph in this study, the gate withholds
\texttt{safe}. Independent event-coverage analysis remains future work.
